% Controllato, teoricamente dovrebbe andare bene sia a lunghezza che a grammatica

\section{Context}\label{sec:context}

The increasing availability of relational data that has been studied in recent years has led to a growing interest in the study of data structures capable of representing complex interactions between entities.
In the literature, such type of relationship were originally managed by means of \textit{graphs}, which are mathematical structures used to model one-to-one relations between objects.
However, many real-world scenarios find this representation too restrictive, when considering interactions between more than just two entities.
\newline
Several real-life scenarios are by default not representable as graphs.
Some examples include conversations between multiple people, usually found online in forums, chat or more generally in social media, while others include the interactions between multiple objects or entities, such as those involved in a bank transaction.
In these types of contexts, the loss of information when representing these interactions as graphs is unavoidable.
\newline
In order to be able to manage this type of data relationship, the concept of \textit{hypergraph} was introduced.
An hypergraph is a generalization of a graph, characterized by the introduction of many-to-many relationships between entities, which are called \textit{hyperedges}.
This type of data structure is particularly useful when the goal is to represent complex interations without losing sensitive information, as it is the case with graphs.
\newline
In real life scenarios, the interactions between entities are not static, but they evolve over time.
New interactions are created, while others tend to disappear, this brings to a complex dynamic system that brings a theoretically static data structure to evolve over time.
This type of challenge makes the management of this temporal evolution a necessity, in order to be able to represent the data always in an accurate way.
\newline
The analysis of large dynamic structures also introduces important computational problems.
In particular, it is often necessary to continuously update the information of interest without being able to recalculate the entire solution from scratch each time.
Consequently, the study of efficient techniques for the maintenance of combinatorial properties and structures in dynamic and temporal environments assumes a central role.